\begin{titlepage}
    \begin{tikzpicture}[remember picture,overlay]
        \node[anchor=north east,inner sep=1cm] at (current page.north east)
        {\includegraphics[width=0.8\textwidth]{../figures/IHF-Logo_platzhalter.png}};
    \end{tikzpicture}
    % \hfill\includegraphics[width=0.5\textwidth]{../figures/IHF-Logo_platzhalter.png}
    \begin{center}
        \vfill

        \Huge
        \textbf{Calibration and Image Reconstruction Techniques
            for FMCW MIMO Radars with Consideration of Antenna Radiation Patterns}
        \vspace{0.5cm}

        \large
        Kalibrations- und Bildrekonstruktionsverfahren für FMCW-MIMO-Radare unter Berücksichtigung der Antennenabstrahlungscharakteristik
        \vspace{1.5cm}

        Master thesis presented by \\
        \textbf{Daniel Gotzens} \\
        Matr.-Nr. 359276

        \vfill

        \normalsize
        Submitted to the Institute of High Frequency Technology \\
        RWTH Aachen University,
        \today\\
        \vfill

        Supervised by: \\
        \large
        Prof. Dr.-Ing. Peter Knott \\
        M. Sc. Benedikt Eliasson \\

        \vspace{1cm}
        In cooperation with \\

        \includegraphics[height=18pt]{../figures/indurad-logo.png}

        indurad GmbH
        Aachen

    \end{center}
\end{titlepage}
\afterpage{\null\newpage}

\subsubsection*{Abstract}
In recent years, FMCW MIMO radars have emerged as highly effective solutions for short-- to medium-range,
high-resolution imaging radar applications. Traditional imaging algorithms rely on the Fast Fourier Transform (FFT),
facilitating real-time processing of radar data to accommodate the high data rates of FMCW MIMO radars.
However, these algorithms exhibit limited flexibility concerning both the shape of the output image and the composition of the input data.
In practical scenarios, this constraint can result in radar image distortion due to near-field effects and antenna gain variations across the array.

This thesis proposes and implements a backprojection-based imaging algorithm for a TDM FMCW MIMO radar operating in the mmWave frequency band.
The thesis evaluates the functionality and stability of the radar sensor, measures its antenna gains, and utilizes them to calibrate the input data.
Subsequently, the implemented algorithm is compared with FFT-based and hybrid approaches.
The hybrid approach combines range estimation using the FFT with direction-of-arrival estimation using backprojection.

The findings reveal that, in comparison to conventional FFT-based algorithms,
backprojection offers the potential for achieving higher image fidelities,
particularly in terms of near-field direction-of-arrival accuracy and signal-to-noise ratio.
